\documentclass[a4j,10pt]{jarticle}
\usepackage[dvipdfm,text={160mm,240mm},centering]{geometry}
\usepackage{newtxtext,newtxmath}
\usepackage[dvipdfmx]{graphicx}
\usepackage{listings,jvlisting}
\usepackage{color}
\usepackage{tabularx}
\usepackage{booktabs}
\usepackage{multirow}

% 担当教員名
\newcommand{\teacherName}{桝井 晃基}
% 提出者名
\newcommand{\studentName}{鶴見 俊介}
% 提出者所属
\newcommand{\studentAff}{基礎工学部 情報科学科3年 計算機科学コース}
% 提出者学籍番号
\newcommand{\studentID}{09B23057}
% 提出者電子メールアドレス
\newcommand{\studentEmail}{u299151f@ecs.osaka-u.ac.jp}
% 課題提出日
\newcommand{\submitted}{\today}
% 課題締め切り日
\newcommand{\deadline}{2025年10月17日}

\lstset{
    basicstyle = {\ttfamily},
    backgroundcolor  = {\color[gray]{.99}},
    breaklines = true,
    columns = [l]{fullflexible},
    commentstyle = {\smallitshape},
    frame = {tb},
    identifierstyle = {\small},
    keepspaces=true,
    keywordstyle = {\small\bfseries},
    lineskip = -0.5ex,
    ndkeywordstyle = {\small},
    numbers = left,
    numbersep = 1zw,
    numberstyle = {\scriptsize},
    stepnumber = 1,
    stringstyle = {\small\ttfamily},
    xleftmargin = 3zw,
    xrightmargin = 0zw
}
\renewcommand{\lstlistingname}{プログラム}

\begin{document}
\begin{titlepage}
\begin{center}
{\Huge\bfseries 情報科学演習D} \\
\vspace{1cm}
{\LARGE\bfseries 課題1}
\vspace{5cm}

{\large
\begin{tabular}{rcl}
担当教員 & : & \teacherName \\
提出者   & : & \studentName \\
所属/学年   & : & \studentAff\\
学籍番号 & : & \studentID \\
電子メール & : & \texttt{\studentEmail}\\
&& \\
提出日 & : & \submitted \\
締切日 & : & \deadline
\end{tabular}
}
\end{center}
\end{titlepage}

\setcounter{tocdepth}{2}
% \tableofcontents
\newpage

% \lstinputlisting[caption=サーバプログラムの接続手順, firstline=1, firstnumber=1, lastline=19, label=prog:server]{codes/server.txt}


\section{システムの仕様}
\section{課題達成の方針と設計}
\section{実装プログラム}
\section{考察や工夫点}
\section{感想}


\end{document}